\section{Background and Motivating Example}
\label{sec:background}

\subsection{Cross-Module Dependencies and Heuristic Blind Spots}
\label{sec:motivating}

In large multi-module repositories, code changes in shared infrastructure or low-level utility libraries frequently induce silent failures in distant consumer modules that share neither common directories nor recent historical failure co-occurrences. 

Consider an illustrative case from Apache ShardingSphere: a developer introduced a single-line modification inside the shared infrastructure module (\texttt{shardingsphere-infra-common}) that invoked \texttt{EXECUTOR\_SERVICE.shutdown()} within the schema-building routine. While syntactically local, this change terminated the shared background thread pool required during database connection initialization. Consequently, four downstream test suites in the database driver module (\texttt{shardingsphere-jdbc-core}) failed during subsequent integration runs.

Standard heuristic selection strategies are fundamentally blind to this interaction:
\begin{itemize}
    \item \textbf{File- and path-based heuristics} fail because the modified infrastructure classes and the failing JDBC test suites share no common subdirectories, packages, or filename stems.
    \item \textbf{History- and telemetry-based guards} skip the tests entirely because the failing JDBC suites had passed consistently on prior builds and exhibited zero recent failure history.
\end{itemize}
Moreover, the diff contains no textual references to JDBC or data sources. To bridge this semantic gap, a test selector must understand the operational preconditions and runtime resources exercised by test suites. By summarizing test suites offline (Section~\ref{sec:indexing}) and computing dense semantic embeddings against the commit diff, our semantic retrieval layer assigns cosine similarities exceeding 0.60 to the affected JDBC suites, correctly triggering their execution without requiring heavyweight runtime bytecode instrumentation.

\subsection{Conformal risk control}
\label{sec:crc}

Conformal risk control calibrates a parameter $\lambda$ on held-out data so that the expected value of a monotone loss on a new point stays below a target $\alpha$~\cite{angelopoulos2024crc}. With a loss bounded by $B$ and $n$ calibration points, it selects $\lambda$ using the empirical risk inflated by a term of order $B/(n+1)$. The result is a statement about expected loss over draws of the calibration data and the new point. It is marginal, so it says nothing about any individual build. An alternative controls the risk with high probability over the calibration set~\cite{bates2021rcps}, which we do not use here.

The result assumes exchangeable data, and CI builds ordered in time are not exchangeable. Barber et al.~\cite{barber2023beyond} bound the coverage degradation for conformal prediction when fixed weights favor observations closer to the test point, in terms of the total variation distance between the data and a version with points swapped. Farinhas et al.~\cite{farinhas2024nonexchangeable} extend this to risk control for bounded monotone losses with data-independent weights. The expected loss is then at most $\alpha$ plus a term proportional to those distances. The term cannot be computed without knowing the distribution, so weighting limits the damage from drift without removing it. Section~\ref{sec:calibration} states which of these properties our method has.
